We present here our main construction of an efficient PPB for threshold comparison.
Let $\pker$ and $\pkem$ be lossy public-key encryption schemes and define the following two relations:
\begin{align*}
	\kgrel^{} = & \left\{
	((pk', \com), (\thr, sk', r_k, r_c)) \::\:
	\begin{gathered}
		(sk', pk') = \kgen'(\lambda, \revpp, \matpp, \base, \numdigits, \thr; r_k) \\
		\wedge \: \com = \commit(\comparams, \thr; r_c)
	\end{gathered}
	\right\}              \\
	\erel^{} =  & \left\{
	((Z', \com, pk), (\uval, \msg, r_e, r_c)) \::\:
	\begin{gathered}
		Z' = \escrow'(\revpp, \matpp, \numdigits, pk, \uval, \msg; r_e) \\
		\wedge \: \com = \commit(\comparams, (\uval, \msg); r_c)
	\end{gathered}
	\right\}
\end{align*}

Let $\kglang^{}$ and $\elang^{}$ be the corresponding languages to these relations, and
$\kgps^{}$, $\eps^{}$ be corresponding NIZKs for these languages.
Our construction $(\maxt, \base)$-TCPPB can then be seen in \cref{fig:ppb-comp-const-01}, \cref{fig:ppb-comp-const-02}, and \cref{fig:ppb-comp-const-03}.
$\maxt$ and $\base$ are assumed to be implicit inputs to all algorithms, and $\msgspace_\rev$ and $\msgspace_\mat$ are the message spaces of $\pke_\rev$ and $\pke_\mat$, respectively, and are both assumed to be finite, cyclic, abelian groups.

\begin{figure}[h]
	\begin{pchstack}[boxed, center, space=1em]

		\procedure[linenumbering]{$\setup(\secparam, \comparams)$}{
			(\crs_{KG}, \td_{KG}) \leftarrow \zksetup_{KG}(\secparam) \\
			(\crs_{E}, \td_{E}) \leftarrow \zksetup_{E}(\secparam) \\
			\revpp \leftarrow \pker.\setup(\secparam) \\
			\matpp \leftarrow \pkem.\setup(\secparam) \\
			\numdigits = \lceil \log_{\base}(\maxt) \rceil \\
			\Lambda = (\lambda, \comparams, \revpp, \matpp, \crs_{KG}, \crs_{E}, \numdigits) \\
			\pcreturn(\ppbpp, (\td_{KG}, \td_{E}))
		}

	\end{pchstack}
	\caption{A $\fthr$-PPB construction $(\maxt, \base)$-TCPPB.}
	\label{fig:ppb-comp-const-01}
\end{figure}

\begin{figure}[ht]
	\begin{pcvstack}[boxed, center, space=1em]


		\begin{pchstack}

			\begin{pcvstack}[space=1em]
				\procedure[linenumbering]{$\kgen\left(\Lambda, \thr, r_{KG}\right)$}{
					\pcparse (\lambda, \comparams, \pkepp^{\rev}, \pcskipln \\
					\t \pkepp^{\mat}, \crs_{KG}, \crs_{E}, \numdigits) = \Lambda \\
					(sk', pk') \xleftarrow{r_k} \pcskipln \\
					\t \kgen'(\lambda, \revpp, \matpp, \numdigits, \thr) \\
					\com_{KG} \leftarrow \commit(\comparams, \thr; r_{KG}) \\
					\zkproof_{KG} \leftarrow \pcskipln \\
					\zkprove_{KG}(\crs_{KG}, (pk', \com_{KG}), \pcskipln \\
					\t (\thr, sk', r_k, r_{KG})) \\
					pk \leftarrow (pk', \com_{KG}, \zkproof_{KG}) \\
					sk \leftarrow (sk', pk) \\
					\pcreturn (sk, pk)
				}

				\procedure[linenumbering]{$\kgen'\left(\lambda, \revpp, \matpp, \numdigits, \thr\right)$}{
					\pcfor i \in [\numdigits]\\
					\pcind \pcfor j \in [0, b] \\
					\t[2] \pcif j \leq t[i] \\
					\t[3] \revpk_i^j \leftarrow \pker.\kgenlossy(\revpp) \\
					\t[3] \matpk_i^j \leftarrow \pkem.\kgenlossy(\matpp) \\
					\t[3] \pairsk_i^j \leftarrow \bot \\
					\t[2] \pcelse \\
					\t[3] \revsk_i^j, \revpk_i^j \leftarrow \pker.\kgen(\revpp) \\
					\t[3] \matsk_i^j, \matpk_i^j \leftarrow \pkem.\kgen(\matpp) \\
					\t[3] \pairsk_i^j \leftarrow (\revsk_i^j, \matsk_i^j)\\
					sk' = \{\pairsk_i^j\}_{i \in [\numdigits], j \in [0 , b - 1]} \\
					pk' = \{\pairpk_i^j\}_{i \in [\numdigits], j \in [0 ,b - 1]} \\
					\pcreturn (sk', pk')
				}

			\end{pcvstack}

			\begin{pcvstack}[space=1em]

				\procedure[linenumbering]{$\escrow(\Lambda, pk, (\uval, \msg), r_{E})$}{
					\pcparse (\lambda, \comparams, \pkepp^{\rev}, \pkepp^{\mat}, \crs_{KG}, \crs_{E}, \numdigits) \pcskipln \\
					\t = \Lambda \\
					\pcparse (pk', \com_{KG}, \zkproof_{KG}) = pk \\
					\pcif \verpk(\ppbpp, pk, \kgcom) = 0 \pcthen \\
					\t \pcreturn \bot \\
					Z' \xleftarrow{r_e} \escrow'(\revpp, \matpp, \numdigits, pk', v, m) \\
					\com_{E} \leftarrow \commit(\comparams, (v, m); r_{E}) \\
					\zkproof_{E} \leftarrow \pcskipln \\
					\t \zkprove_{E}(\crs_{E}, (Z', \com_{E}, pk'), (v, m, r_e, \er)) \\
					\pcreturn (Z', \zkproof_{E})
				}

				\procedure[linenumbering]{$\escrow'(\revpp, \matpp, \numdigits, pk, v, m)$}{
					\otpk^* \sample \msgspace_{\rev} \\
					\otpct^* := \otpk^* \cdot \msg \\
					\otpk_1,\dots,\otpk_{\numdigits - 1} \sample \msgspace_{\mat} \\
					\otpk_0 := \msgspace_{\mat}.e \\
					\pcfor i \in [n] \\
					\t \revct_i \leftarrow \pker.\enc(\revpp, \revpk_i^{\uval[i]}, \otpk^*) \\
					\t \otpct_i^{\rev} := \otpk_{i - 1} \cdot \revct_i \\
					\t \pcif i < \numdigits \\
					\t[2] \otpct^{\mat}_i := \otpk_{i - 1}\cdot \otpk_i \\
					\t[2] \matct_i \leftarrow \pkem.\enc( \pcskipln \\
					\t[4] \matpp, \matpk_i^{\uval[i] + 1}, \otpct^{\mat}_i) \\
					\t \pcelse\\
					\t[2] \matct_i \leftarrow \bot \\
					\t \pairct_i := (\otpct_i^{\rev}, \matct_i) \\
					\pcreturn (\{\pairct_i\}_{i \in [\numdigits]}, \otpct^*)
				}
			\end{pcvstack}
		\end{pchstack}

	\end{pcvstack}
	\caption{$(\maxt, \base)$-TCPPB continued (key generation and escrow).}
	\label{fig:ppb-comp-const-02}
\end{figure}


\begin{figure}[ht]
	\begin{pchstack}[boxed, center, space=1em]
		\begin{pcvstack}[space=1em]
			\procedure[linenumbering]{$\verpk(\Lambda, pk, \com_{KG})$}{
				\pcparse (\lambda, \comparams, \crs_{KG}, \crs_{E},  \pkepp, \maxt, \base, \numdigits) = \Lambda \\
				\pcparse (pk', \com_{KG}', \zkproof_{KG}) = pk \\
				\pcreturn \zkverify_{KG}(\crs_{KG}, \zkproof_{KG}, (pk', \com_{KG})) \pcskipln \\
				\t \wedge (\com_{KG}'  = \com_{KG})
			}
			\procedure[linenumbering]{$\verescrow(\Lambda, pk, \com_{E}, Z = (Z', \zkproof_{E}))$}{
				\pcparse (\lambda, \comparams, \crs_{KG}, \crs_{E},  \pkepp, \maxt, \base, \numdigits) = \Lambda \\
				\pcparse (pk', \com_{KG}', \zkproof_{KG}) = pk \\
				\pcreturn \verpk(\Lambda, pk, \com_{KG}) \pcskipln \\
				\t \wedge \: \zkverify_{E}((Z', \com_{E}, pk'), \zkproof_{E})
			}
			\procedure[linenumbering]{$\dec(\Lambda, sk, \com_{E}, Z = (Z', \zkproof_E))$}{
				\pcparse (\lambda, \comparams, \pkepp^{\rev}, \pkepp^{\mat}, \crs_{KG}, \crs_{E}, \numdigits) = \Lambda \\
				\pcparse (\{(\revsk_i^j, \matsk_i^j)\}_{i \in [\numdigits], j \in [0 ,b]}, pk) \pcskipln \\
				\t = sk \\
				\pcif \verescrow(\Lambda, pk, \com_{E}, Z) = 0 \:\pcreturn \bot \\
				\pcparse (\{\pairct_i\}_{i \in [\numdigits]}, \otpct^*) = Z' \\
				\tilde{\otpk}_0 := \msgspace_{\mat}.e \\
				\pcfor i \in [n] \\
				\t \pcparse \pairct_i  = (\otpct^{\rev}_i, \matct_i)\\
				\t \revct_i := \tilde{\otpk}_{i - 1}^{-1} \cdot \otpct^{\rev}_i \\
				\t \pcfor j \in [\thr[i] + 1, \base] \\
				\t[2] \tilde{\otpk}^* := \pker.\dec(\revpp, \revsk_i^j, \revct_i) \\
				\t[2] \pcif \tilde{\otpk}^* \neq \bot \text{ } \pcreturn (\tilde{\otpk}^*)^{-1} \cdot \otpct^* \\
				\t \pcif i < n \\
				\t\t \otpct^{\mat}_{i} := \pkem.\dec(\matpp, \matsk_i^{\thr[i] + 1}, \matct_i) \\
				\t\t \tilde{\otpk}_i := \tilde{\otpk}_{i - 1}^{-1}\cdot \otpct^{\mat}_i \\
				\pcreturn \bot
			}
		\end{pcvstack}
	\end{pchstack}
	\caption{$(\maxt, \base)$-TCPPB continued (verification and decryption).}
	\label{fig:ppb-comp-const-03}
\end{figure}



We can now state our main theorem.

\begin{theorem}
	\label{thm:ppbs-main}
	If $\pker$ is a weakly robust, lossy, \wlap{} public-key encryption scheme with message space $\msgspace_\rev$ and ciphertext space $\msgspace_\mat$, $\pkem$ is a lossy, $\mathcal{U}_{\msgspace_\mat}$-\slap{} public-key encryption scheme with message space $\msgspace_\mat$, $\msgspace_\rev$ and $\msgspace_\mat$ are both finite, cyclic, abelian groups, $(\pker, \pkem)$ satisfy lossy correlation resistance, $\comscheme = (\csetup, \commit)$ is a secure commitment scheme, $\kgps$ is zero-knowledge and strong simulation extractable, and $\eps$ is zero-knowledge and satisfies knowledge soundness, then \schemename{} is a secure $\fthr$-PPB scheme.
\end{theorem}


\FloatBarrier
\subsection{Proofs}
As our construction uses the generic commit-and-prove technique from \cite{EC:KohLysNgu23}, the proofs of correctness of $\verescrow$ and $\verpk$, follow from the same arguments, and we omit the proofs here.

\begin{lemma}
	\label{lemma:main-soundness}
	If $\Pi_{E}$ is a knowledge-sound NIZK, $\comscheme$ is a binding commitment scheme, $\pker$ satisfies weak-robustness, and $\pker$, $\pkem$ are correct and satisfy lossy correlation resistance, then \schemename{} satisfies soundness.
\end{lemma}

\begin{proof}
	Suppose for the sake of contradiction that there exists an adversary $\adv$ such that $\advantage{\soundgame}{\mathsf{TCPPB}, \adv} = \nu(\secpar)$ is non-negligible.
	In the adversary's first stage it outputs $(\thr, \kgr)$, and the challenger sets $(sk, pk) \leftarrow \kgen(\ppbpp, \thr, \kgr)$.
	In its second stage, the adversary $\adv$ outputs $(\ecom, (\uval, \msg), \er, \esc = (\esc', \zkproof_{E}))$.
	We can divide our analysis into the following three cases depending on the properties of these outputs.
	\begin{enumerate}
		\item There exist some $\er'$ and $(\uval', \msg') \neq (\uval, \msg)$ such that $\commit((\uval, \msg), \er) = \commit((\uval', \msg'), \er')$ and $\esc = \escrow(\ppbpp, pk, (\uval', \msg'), \er')$.
		\item $Z' \neq \escrow'(\pkerpp, \pkempp, \numdigits, pk, \uval, \msg)$ and (1) does not hold.
		\item Neither (1) nor (2) hold (i.e. $\adv$'s output is well-formed).
	\end{enumerate}
	Let the probability of $\adv$'s outputs satisfying property $j$ be denoted by $\nu_j(\secpar)$.
	If $\nu(\lambda)$ is non-negligible, then at least one of $\nu_1(\secpar)$, $\nu_2(\secpar)$, $\nu_3(\secpar)$ must be non-negligible.

	${}$

	Suppose that $\nu_2(\secpar)$ is non-negligible.
	In this case $\adv$ produced a proof of a false statement and we can construct a reduction $\bdv$ to the knowledge-soundness of $\Pi_{E}$.
	$\bdv$ runs $\adv$, simulates the generation of $pk$, receives $(\ecom, (\uval, \msg), \er, \esc = (\esc', \zkproof_{E}))$, and outputs $((Z', \ecom, pk), \zkproof_{E})$.
	As the statement is false any extractor must fail, meaning $\prob{\bdv \text{ wins}}$ is non-negligible, thus contradicting our assumption that $\Pi_{E}$ satisfies knowledge soundness.
	Therefore, $\nu_2(\secpar)$ is negligible.

	As $\nu_2(\secpar)$ is negligible we can assume there exists an extractor $\Ext_{\adv}$ that given $\trans_{\adv}$ outputs $(\uval', \msg', r_e, \er)$ such that $\ecom = \commit(\compp, (\uval', \msg'); \er)$ and $\esc'\allowbreak =\allowbreak \escrow'(\allowbreak\pkerpp,\allowbreak \pkempp,\allowbreak \numdigits,\allowbreak pk,\allowbreak \uval',\allowbreak \msg'; r_e)$.

	${}$

	Suppose $\nu_1(\secpar)$ is non-negligible.
	In this case we can break the computational binding property of the commitment scheme by using a reduction that uses the extractor and outputs $(\msg, \er, \msg', r_c)$.
	Therefore, $\nu_1(\secpar)$ must be negligible.

	${}$

	Suppose $\nu_3(\secpar)$ is non-negligible.
	In this case we can break a property of $\pker$ or $\pkem$.
	In all cases, the reduction begins by using the extractor to learn $(\uval', \msg', r_e, \er)$.
	The reduction can then use $r_e$ and $r_k$ to infer the randomness used to generate each individual value in $\escrow'$.
	Let $\rer_i$ and $\mer_i$ be the randomness used at iteration $i$ of $\escrow'$ for $\pker.\enc$ and $\pkem.\enc$, respectively.

	Let $\nu_3^{>}(\lambda)$ denote the probability that when $\nu_3(\secpar)$ is true,  $\uval' > \thr$.
	Let  $\nu_3^{<}(\lambda)$ be the case that $\uval' < \thr$, and $\nu_3^{=}(\secpar)$ be the case that $\uval' = \thr$.
	If $\nu_3(\lambda)$ is non-negligible, then at least one of $\nu_3^{>}(\lambda)$, $\nu_3^{<}(\lambda)$, $\nu_3^{=}(\secpar)$ must be non-negligible.

	${}$

	Suppose $\nu_3^{>}(\lambda)$ is non-negligible.
	In this case, let $d$ be the index for which $\uval'[d] > \thr[d]$ and $\forall i < d,\: \uval'[i] = \thr[i]$.
	We can split this into further cases depending on the behavior of $\dec$.
	To do so, we define the following notation.

	\newcommand{\badprobmat}[1]{\ensuremath{\nu_{3, #1}^{>,\mat}(\secpar)}}
	\newcommand{\badprobotpM}[1]{\ensuremath{\nu_{3, #1}^{>, \mathsf{OM}}(\secpar)}}
	\newcommand{\badprobotpR}[1]{\ensuremath{\nu_{3, #1}^{>, \mathsf{OR}}(\secpar)}}
	\newcommand{\badprobrev}[2]{\ensuremath{\nu_{3, #1, #2}^{>, \mathsf{\rev}}(\secpar)}}
	\newcommand{\badprobrevF}{\ensuremath{\nu_{3, d, \uval'[d]}^{>, \mathsf{*}}(\secpar)}}
	\begin{itemize}
		\item Let \badprobmat{i} denote the probability that $\pkem.\dec(\pkempp, \matsk_i^{\thr[i] + 1}, \mct_i) \allowbreak\neq\allowbreak \otpk_{i - 1} \cdot \otpk_i$ for some $i \in [d - 1]$.
		\item Let \badprobotpM{i} denote the probability that $\tilde{\otpk}_i \neq \otpk_i$ for some $i \in [0, d - 1]$
		\item Let \badprobotpR{i} denote the probability that $\revct_i \neq \pker.\enc(\pkerpp, \revpk_i^{\uval'[i]}, \otpk^*; \rer_i)$ for some $i \in [d]$.
		\item Let \badprobrev{i}{j} denote the probability that $\pker.\dec(\pkerpp, \revsk_i^j, \revct_i) \neq \bot$ for some $i \in [d - 1]$ and $j \in [\thr[i] + 1, \base]$, or for $i = d$ and $j \in [\thr[i] + 1, \uval'[d] - 1]$.
		\item Let \badprobrevF{} denote the probability that $\pker.\dec(\pkerpp, \revsk_d^{\uval'[d]}, \revct_d) \neq \otpk^*$.
	\end{itemize}

\end{proof}
